% ===== 第 5 章 §5.4 前段（印刷 pp.203–217 / PDF p.218–232）· tr-01 =====
% 边界：c05b 结束于 Proposition 5.28 的证明中途（proof 环境未闭合，其最后一句译到 p.203 首行的前一半）。
% 本块块首不开 proof 环境，直接续写；在 p.203 证明收尾处补上 proof 环境的结束（□ 由 amsthm 自动给出）。
% §5.4 的 \subsection 标题在 c05b；本块含 \subsubsection 5.4.4 / 5.4.5 / 5.4.6（自动编号）。
% tag 区间 5.10–5.17：p.209 (5.10)(5.11)(5.12)、p.210 (5.13)、p.211 (5.14)、p.212 (5.15)(5.16)、p.216 (5.17)。
% 断言 5.36 / 5.37 用 preamble 的 claim 环境（Lead 已加入 \newtheorem{claim}{断言}，与原文 “Claim” 完全对应）。

% [原文 p.203]
另一方面，若 $I(h_1)\cap I(h_2)=\varnothing$，则 $h_1$ 与 $h_2$ 可以粘接得到一棵新树 $h$，而
$\mathfrak D_P$ 由三个元素 $\mathfrak d_1$、$\mathfrak d_2$ 与 $\mathfrak d$ 组成，分别对应于 $h_1$、$h_2$
与 $h$。此时
\[
f_{\mathfrak d_i} = 1 + w_\Gamma(E_{\mathrm{out},i})\operatorname{Mult}(h_i)z^{\Delta(h_i)}u_{I(h_i)}
\]
对 $i=1,2$，并且
\begin{align*}
f_{\mathfrak d} &= 1+w_\Gamma(E_{\mathrm{out}})\operatorname{Mult}(h)z^{\Delta(h)}u_{I(h)}\\
&= 1+w_\Gamma(E_{\mathrm{out}})\operatorname{Mult}(h_1)\operatorname{Mult}(h_2)\operatorname{Mult}_V(h)
  z^{\Delta(h_1)+\Delta(h_2)}u_{I(h_1)}u_{I(h_2)}\\
&= 1+w_\Gamma(E_{\mathrm{out}})\operatorname{Mult}(h_1)\operatorname{Mult}(h_2)
  w_{\Gamma_1}(E_{\mathrm{out},1})w_{\Gamma_2}(E_{\mathrm{out},2})\cdot\\
&\qquad\cdot |m'_1\wedge m'_2|z^{\Delta(h_1)+\Delta(h_2)}u_{I(h_1)}u_{I(h_2)}\\
&= 1+w_\Gamma(E_{\mathrm{out}})\operatorname{Mult}(h_1)\operatorname{Mult}(h_2)\cdot\\
&\qquad\cdot |r(\Delta(h_1))\wedge r(\Delta(h_2))|z^{\Delta(h_1)+\Delta(h_2)}u_{I(h_1)}u_{I(h_2)}.
\end{align*}
这里我们用到了如下事实：若把平衡条件对 $\Gamma'_i$ 的所有顶点求和，则得到
\[
w_{\Gamma'_i}(E_{\mathrm{out},i})m'_i = r(\Delta(h_i)).
\]
于是由例 5.24，$\theta_{\gamma_0,\mathfrak D_P}$ 是恒等映射。

对一般情形，$\mathfrak D_P$ 中有某有限多条直线，以及若干射线。设 $\mathfrak D_P$ 中有三条直线分别对应于树
$h_1,h_2,h_3$，且 $I(h_1),I(h_2),I(h_3)$ 两两不交。则如上面两条直线的情形那样，这些树可以在 $P$ 处粘接，
得到一棵带四价顶点的 Maslov 指数零树。然而由于 $P_1,\dots,P_k$ 处于一般位置，没有 Maslov 指数零树含价
$>3$ 的顶点。因此这种情形不会发生。另一方面，给定两条直线对应于树 $h_1,h_2$ 且
$I(h_1)\cap I(h_2)=\varnothing$，这两棵树可以如上在 $P$ 处粘接，得到一棵新的 Maslov 指数零树。于是
$\mathfrak D_P$ 中的射线与满足 $I(h_1)\cap I(h_2)=\varnothing$ 的、对应于树 $h_1,h_2$ 的直线对
$\mathfrak d_1,\mathfrak d_2\in\mathfrak D_P$ 一一对应。因此可以写成
\[
\mathfrak D_P = \{\mathfrak d_1,\dots,\mathfrak d_n\}\cup\bigcup_{j=1}^m \mathfrak D_i
\]
其中 $\mathfrak d_1,\dots,\mathfrak d_n$ 是对应于这样的树 $h$ 的直线：对任何出射边经过 $P$ 的 Maslov 指数零树
$h'$ 都有 $I(h)\cap I(h')\ne\varnothing$；而 $\mathfrak D_1,\dots,\mathfrak D_m$ 是散射图，每个由两条直线与
一条射线组成，其中直线对应于满足 $I(h_1)\cap I(h_2)=\varnothing$ 的树 $h_1$ 与 $h_2$，射线对应于把 $h_1$ 与
$h_2$ 在 $P$ 处粘接所得的树。

% [Tutor Note: 原文 p.203，$\mathfrak D_P=\{\mathfrak d_1,\dots,\mathfrak d_n\}\cup\bigcup_{j=1}^m\mathfrak D_i$
%  的并集指标为 $j$，被并的对象却写成 $\mathfrak D_i$——这是原书的指标笔误，应为 $\mathfrak D_j$（紧接的
%  $\prod_{j=1}^m\theta_{\gamma_0,\mathfrak D_j}$ 也证实这一点）。译文照录不改。]

现在计算 $\theta_{\gamma_0,\mathfrak D_P}$ 是一道交换子练习。注意若 $\mathfrak d_1,\mathfrak d_2\in\mathfrak D_P$
对应于两棵树 $h_1,h_2$ 且 $I(h_1)\cap I(h_2)\ne\varnothing$，则如前所述，$\theta_{\gamma_0,\mathfrak d_1}$ 与
$\theta_{\gamma_0,\mathfrak d_2}$ 交换。于是利用这一交换性，可以写成
\[
\theta_{\gamma_0,\mathfrak D_P}
= \left(\prod_{i=1}^n \theta_{\gamma_0,\mathfrak d_i}\circ\theta^{-1}_{\gamma_0,\mathfrak d_i}\right)
\circ\prod_{j=1}^m\theta_{\gamma_0,\mathfrak D_j}.
\]
当然 $\theta_{\gamma_0,\mathfrak d_i}\circ\theta^{-1}_{\gamma_0,\mathfrak d_i} = \mathrm{Id}$，而
$\theta_{\gamma_0,\mathfrak D_j} = \mathrm{Id}$（由已处理的特殊情形）。因此在这一般情形下
$\theta_{\gamma_0,\mathfrak D_P} = \mathrm{Id}$。
\end{proof}

% [原文 p.204]
\begin{remark}
注意 $\mathfrak D = \mathfrak D(\Sigma,P_1,\dots,P_k)$ 中以 $P_i$ 为端点的射线，与
$(X_\Sigma,P_1,\dots,P_{i-1},P_{i+1},\dots,P_k)$ 中以 $P_i$ 为边界的 Maslov 指数二圆盘一一对应。

事实上，任取这样一个 Maslov 指数二圆盘，延拓出射边得到一棵热带树，我们便可在这条出射边上标出映到 $P_i$ 的点，
从而得到 $(X_\Sigma,P_1,\dots,P_k)$ 中 Maslov 指数零的热带树。$\mathfrak D$ 中对应的射线以 $P_i$ 为端点。反之，
给定 $\mathfrak D$ 中以 $P_i$ 为端点的射线，它对应于一棵 Maslov 指数零树，使得与 $E_{\mathrm{out}}$ 相邻的
顶点 $V$ 是某条映到 $P_i$ 的标记边 $E_x$ 的顶点。在 $V$ 处切开这棵树并去掉映到 $P_i$ 的标记边，就得到以
$P_i$ 为边界的 Maslov 指数二圆盘。

此外，由 $P_j$ 的一般位置，$\mathfrak D$ 中不含这样的射线：它含有点 $P_i$，却不以 $P_i$ 为端点。
\end{remark}

\subsubsection{断线}
这种散射图方法的一个好处是，借助我们所谓的\emph{断线}，可以容易地描述以一般点 $Q$ 为边界的 Maslov 指数二圆盘：

% [原文 p.205]
\begin{definition}
\emph{断线}是指连续、逆紧的分片线性映射
\[
\beta : (-\infty,0] \to M_\R
\]
其端点 $Q=\beta(0)$，并带有如下所述的一些附加数据。设
\[
-\infty = t_0 < t_1 < \dots < t_n = 0
\]
是使 $\beta|_{(t_{i-1},t_i)}$ 为线性的最小实数集。则对每个 $1\le i\le n$，给定单项式
$c_iz^{m^\beta_i}\in \C[T_\Sigma]\otimes_\C R_k[\![y_0]\!]$ 作为附加数据，其中
$m^\beta_i\in T_\Sigma\setminus K_\Sigma$ 且 $0\ne c_i\in R_k$。此外这些数据应满足下列性质：
\begin{enumerate}[label=(\arabic*),leftmargin=*]
\item 对每个 $i$，$r(m^\beta_i) = -\beta'(t)$，其中 $t\in(t_{i-1},t_i)$。
\item $m^\beta_1 = t_\rho$ 对某个 $\rho\in\Sigma^{[1]}$ 成立，且 $c_1 = 1$。
\item $\beta(t_i)\in\operatorname{Supp}(\mathfrak D(\Sigma,P_1,\dots,P_k))\setminus
\operatorname{Sing}(\mathfrak D(\Sigma,P_1,\dots,P_k))$，$1\le i\le n$。
\item 若 $\beta(t_i)\in\mathfrak d_1\cap\dots\cap\mathfrak d_s$（这个交必是一维的），则
$c_{i+1}z^{m^\beta_{i+1}}$ 是
\[
(\theta_{\beta,\mathfrak d_1}\circ\dots\circ\theta_{\beta,\mathfrak d_s})(c_iz^{m^\beta_i})
\]
中的一项。
\end{enumerate}
这里我们指的是如下意思。设 $f_{\mathfrak d_j} = 1+c_{\mathfrak d_j}z^{m_{\mathfrak d_j}}$，$1\le j\le s$，且
$c^2_{\mathfrak d_j}=0$；设 $n\in N$ 本原、与所有 $\mathfrak d_j$ 正交，并选取得使
\begin{align*}
(\theta_{\beta,\mathfrak d_1}\circ\dots\circ\theta_{\beta,\mathfrak d_s})(c_iz^{m^\beta_i})
&= c_iz^{m^\beta_i}\prod_{j=1}^s(1+c_{\mathfrak d_j}z^{m_{\mathfrak d_j}})^{\langle n,r(m^\beta_i)\rangle}\\
&= c_iz^{m^\beta_i}\prod_{j=1}^s\bigl(1+\langle n,r(m^\beta_i)\rangle c_{\mathfrak d_j}z^{m_{\mathfrak d_j}}\bigr).
\end{align*}
则必须有
\[
c_{i+1}z^{m^\beta_{i+1}} = c_iz^{m^\beta_i}\prod_{j\in J}\langle n,r(m^\beta_i)\rangle c_{\mathfrak d_j}z^{m_{\mathfrak d_j}}
\]
对某个指标集 $J\subseteq\{1,\dots,s\}$。我们把这理解为 $\beta$ 在时刻 $t_i$ 被射线
$\{\mathfrak d_j \mid j\in J\}$ 所\emph{折弯}。
\end{definition}

% [Tutor Note: 原文 p.204–205，定义 5.30。三点易错处：① 单项式取在
%  $\C[T_\Sigma]\otimes_\C R_k[\![y_0]\!]$ 中，其中 $R_k[\![y_0]\!]$ 是形式幂级数环（书中用空心方括号），
%  $y_0$ 是与散射图初始射线方向对应的形式变量。② $m^\beta_i\in T_\Sigma\setminus K_\Sigma$：排除“落在主锥
%  （度数锥）里的指数”，这正是把计数写成 $\C[T_\Sigma]$ 中元素、从而与量子上同调对接的关键。
%  ③ 条件 (1) 的负号 $r(m^\beta_i)=-\beta'(t)$：$\beta$ 随 $t$ 增大走向终点，而 $m^\beta_i$ 记录该段的单项式指数，
%  二者相差一个 $r$ 前的负号；核对锚点：当 $\beta$ 是直线 $\beta(t)=Q+t\cdot m$（$t\le0$）时，$\beta'(t)=m$，
%  而指向终点的向量是 $m$ 而非 $-m$，故条件写作 $r(m^\beta_i)=-\beta'(t)$。]

\begin{figure}[htbp]
\centering
\includegraphics[width=0.88\textwidth]{f5_12.png}
\figtitle{12}{以 $Q$ 为终点的各条断线。}
\end{figure}

% [Tutor Note: 原文 p.205，图 12。图中每条折线的各线段都标以相应单项式（如 $1+x_0x_1u_1u_2$），
%  奇点 $P_1,P_2$ 是散射图 $\mathfrak D$ 中射线的端点；虚线是 $\mathfrak D$ 中的直线（含重叠的平行直线）。
%  注意各单项式的指数中 $x_0,x_1,x_2$ 是 $T_\Sigma$ 的基、$u_1,u_2$ 只出现在指数 $\ge1$ 处（对应 $u_{I(h)}$）。]

\begin{example}
仍在 $\PP^2$、$k=2$ 的情形，图 12 画出了以给定的点 $Q$ 为 $\beta(0)$ 的那些断线。每条断线的各线段都标以
相应的单项式。
\end{example}

\begin{proposition}
若 $Q\notin\operatorname{Supp}(\mathfrak D(\Sigma,P_1,\dots,P_k))$ 是一般点，则以 $Q$ 为终点的断线与以 $Q$ 为
边界的 Maslov 指数二圆盘一一对应。此外，若 $\beta$ 是对应于圆盘 $h$ 的断线，而 $cz^m$ 是联系于 $\beta$ 的最后
一段（即端点为 $t_n=0$ 的那一段）的单项式，则
\[
cz^m = \operatorname{Mono}(h).
\]
\end{proposition}

\begin{proof}
我们先证明下面的断言：

\para{断言。}设 $h : \Gamma' \to M_\R$ 是 $(X_\Sigma,P_1,\dots,P_k)$ 中以 $Q'\in M_\R$ 为边界的 Maslov 指数
二圆盘。此外设 $\Gamma'$ 除 $V_{\mathrm{out}}$ 以外的所有顶点都是三价的，且 $h$ 不能在以 $Q'$ 为边界的
Maslov 指数二圆盘族中连续形变。则存在唯一确定的子集 $\Xi = \Xi(h)\subseteq\overline\Gamma$，它是 $\overline\Gamma$ 的边的并、
同胚于 $[-\infty,0]$，把 $\overline\Gamma^{[0]}_\infty\setminus\{V_{\mathrm{out}}\}$ 中的某点连接到 $V_{\mathrm{out}}$，
并满足：
\begin{enumerate}[label=(\arabic*),leftmargin=*]
\item $\Xi$ 与所有 $\partial E_{p_i}$ 不相交。
\item $h$ 在 $\Gamma'\setminus\Xi$ 的任一连通分支的闭包上的限制都是 Maslov 指数零圆盘。
\end{enumerate}

\begin{proof}
我们对 $\Gamma'$ 的顶点数作归纳。若 $\Gamma'$ 只有一个顶点 $V_{\mathrm{out}}$，则 $\Gamma'$ 只有一条边且没有
标记边。此时取 $\Xi(h)$ 为这条边即可。

对归纳步，设 $\Gamma'$ 的出射边为 $E_{\mathrm{out}}$，其端点为 $V_{\mathrm{out}}$ 与 $V$，且
$h(V_{\mathrm{out}}) = Q'$。

% [原文 p.206]
首先我们证明 $V$ 不可能是某条标记边 $E_{p_i}$ 的顶点。否则有 $h(V) = h(p_i) = P_j$ 对某个 $j$ 成立，那么我们
可在 $V$ 处切开 $h$ 并去掉标记边 $E_{p_i}$。这给出一个以 $P_j$ 为边界、但标记点比 $h$ 少一个的圆盘 $h'$。于是
$h'$ 的 Maslov 指数为四，从而由引理 5.6，$h'$ 可以在保持端点 $P_j$ 不动的单参数族中形变。注意对 $h'$ 的小形变，
边 $h'(E'_{\mathrm{out}})$ 的切方向不变。因此 $h'$ 的形变可以延拓为 $h$ 的形变。这与“$h$ 不能形变”矛盾。
故 $V$ 不可能是某条 $E_{p_i}$ 的顶点。

现在在 $V$ 处切开 $h$，令 $\Gamma'_1$ 与 $\Gamma'_2$ 是 $\Gamma'\setminus\{V\}$ 的两个不含 $V_{\mathrm{out}}$ 的
连通分支的闭包。令 $h_i = h|_{\Gamma'_i}$。这给出两个以 $h(V)$ 为边界的圆盘 $h_1,h_2$。我们有
$\mathrm{MI}(h) = \mathrm{MI}(h_1)+\mathrm{MI}(h_2)$。若 $\mathrm{MI}(h_1)\ge4$，则 $h_1$ 可以在保持端点 $h(V)$ 不动的情形下形变，再把这样的形变
与 $h_2$ 粘接，就得到 $h$ 的形变，又是矛盾。于是 $\mathrm{MI}(h_1),\mathrm{MI}(h_2)\le2$，从而必有 $\mathrm{MI}(h_1)=0$ 且
$\mathrm{MI}(h_2)=2$，或反之。

不失一般性，设 $\mathrm{MI}(h_2)=2$。注意 $h_2$ 现在是以 $Q'' = h_2(V)$ 为边界的 Maslov 指数二圆盘。若 $h_2$ 可以在以
$Q''$ 为边界的圆盘族中形变，则把这些形变与 $h_1$ 粘接，就得到 $h$ 的形变，矛盾。于是 $h_2$ 满足断言的假设，
而 $\Gamma'_2$ 的顶点数比 $\Gamma'$ 少。所以由归纳假设 $\Xi(h_2)$ 存在，我们可以取
\[
\Xi(h) = E_{\mathrm{out}}\cup\Xi(h_2).
\]
$\Xi(h)$ 满足所要求的两条性质：因为 $\Xi(h_2)$ 满足，且 $E_{\mathrm{out}}$ 与所有 $\partial E_{p_i}$ 不相交，
并且 $h_1$ 是 Maslov 指数零圆盘。
\end{proof}

% [Tutor Note: 原文 p.205–207，命题 5.32 的证明。核心是断言给出的子集 $\Xi(h)$：它是 $\Gamma'$ 中一条从某个
%  无界顶点连到 $V_{\mathrm{out}}$ 的链，$h|_\Xi$ 就是断线 $\beta$，而 $\Gamma'\setminus\Xi$ 的各连通分支正好给出
%  $\beta$ 各段所附的单项式 $\operatorname{Mono}(h_i)$。断言对 $\Gamma'$ 的顶点数作归纳，其中“$V$ 不可能是任何
%  $E_{p_i}$ 的顶点”一步用的是“Maslov 指数四的圆盘可形变”（引理 5.6）——这是全证明中唯一用到“一般位置”的地方。]

现在固定以 $Q$ 为边界的 Maslov 指数二圆盘 $h : \Gamma'\to M_\R$。由 $Q,P_1,\dots,P_k$ 的一般性，$h$ 满足断言
的假设。取 $\beta = h|_{\Xi(h)}$，我们看到 $\beta$ 是分片线性的。按断线的定义选取
$-\infty = t_0<\dots<t_n = 0$。每个 $t_i$ 对应于 $\Gamma'$ 的一个顶点 $V_i$。当然对 $i\ne n$，
$\Gamma'\setminus\{V_i\}$ 有两个不含 $V_{\mathrm{out}}$ 的连通分支，而断言的证明表明：把 $h$ 限制在其中之一的
闭包上，给出以 $h(V_i)$ 为边界的 Maslov 指数二圆盘，我们现在把它记作 $h_i$。另一个分支同样给出 Maslov 指数零
圆盘。因此 $\beta(t_i)\in\operatorname{Supp}(\mathfrak D(\Sigma,P_1,\dots,P_k))$，$1\le i\le n$。我们取单项式
$c_iz^{m_i}$ 为 $\operatorname{Mono}(h_i)$，并需验证带这些数据的 $\beta$ 现在是断线。条件 (3) 刚刚已证，
条件 (2) 显然。因此对足够一般的 $Q$，只需再验证条件 (1) 与 (4)。我们需要说明：联系于连接 $V_i$ 与
$V_{i+1}$ 的边的单项式 $\operatorname{Mono}(h_{i+1})$ 由联系于连接 $V_{i-1}$ 与 $V_i$ 的边的单项式
$\operatorname{Mono}(h_i)$ 按条件 (4) 产生。设 $\Gamma'$ 中以 $V_i$ 为根的两棵子树为 $g$ 与 $h_i$，且
$\mathrm{MI}(g)=0$。则
\[
\operatorname{Mono}(h_i) = \operatorname{Mult}(h_i)z^{\Delta(h_i)}u_{I(h_i)},
\]
而若 $\mathfrak d$ 是对应于由 $g$ 所得热带树的射线，则
\[
f_{\mathfrak d} = 1+w_{\mathrm{out}}(g)\operatorname{Mult}(g)z^{\Delta(g)}u_{I(g)}.
\]
这里 $w_{\mathrm{out}}(g)$ 表示 $g$ 的出射边的权。把平衡条件在两个圆盘中任一个的每个顶点处求和，可以写出
$-w_{\mathrm{out}}(g)m'_1 = r(\Delta(g))$ 与 $-w_{\mathrm{out}}(h_i)m'_2 = r(\Delta(h_i))$，其中 $m'_i\in M$ 本原，
$m'_1$ 与

% [原文 p.207]
$g$ 的出射边相切、$m'_2$ 与 $h_i$ 的出射边相切，且都指向 $h(V_i)$。特别地，在 $\Xi(h)$ 的参数化选取正确的前提
下，这就给出条件 (1)。

现在可以如下定义自同构 $\theta_{\beta,\mathfrak d}$：取 $n_0 = \pm X_{m'_1}$，符号的选取由要求
$\langle n_0,-m'_2\rangle>0$ 给出；因为关于 $n_0$ 的约定说 $n_0$ 在我们穿过 $\mathfrak d$ 的方向所指的向量上
应为负，而 $m'_2$ 正是这样的向量，故 $n_0$ 在 $-m'_2$ 上为正。于是
$\langle n_0,-m'_2\rangle = |m'_1\wedge m'_2|$，并且
\begin{align*}
\theta_{\beta,\mathfrak d}(\operatorname{Mono}(h_i))
&= \operatorname{Mult}(h_i)z^{\Delta(h_i)}u_{I(h_i)}\\
&\quad + \operatorname{Mult}(h_i)\operatorname{Mult}(g)|m'_1\wedge m'_2|
 w_{\mathrm{out}}(g)w_{\mathrm{out}}(h_i)z^{\Delta(h_i)+\Delta(g)}u_{I(h_i)}u_{I(g)}.
\end{align*}
我们关心的正是第二项，它等于
\begin{align*}
&\operatorname{Mult}(h_i)\operatorname{Mult}(g)\operatorname{Mult}_{V_i}(h_{i+1})
 z^{\Delta(h_{i+1})}u_{I(h_{i+1})}\\
&\qquad= \operatorname{Mult}(h_{i+1})z^{\Delta(h_{i+1})}u_{I(h_{i+1})}
 = \operatorname{Mono}(h_{i+1}),
\end{align*}
如所求。

反之，给定断线 $\beta$，容易构造相应的 Maslov 指数二圆盘：在 $\beta$ 的定义域 $(-\infty,0]$ 上，于每个折弯点
处粘接 Maslov 指数零圆盘。特别地，若 $\beta(t_i)$ 落在射线
$\mathfrak d_1,\dots,\mathfrak d_s\in\mathfrak D(\Sigma,P_1,\dots,P_k)$ 中，且 $\beta$ 在时刻 $t_i$ 被这些
射线中的子集 $\{\mathfrak d_j\mid j\in J\}$ 折弯，则对每个 $j\in J$，把以 $\beta(t_i)$ 为端点、对应于
$\mathfrak d_j$ 的 Maslov 指数零圆盘粘接到 $t_i\in(-\infty,0]$ 上。（注意由 $P_i$ 与 $Q$ 的一般位置，事实上
可以假定 $\#J = 1$。）显然这逆转了上面从 Maslov 指数二圆盘得到断线的过程。
\end{proof}
\subsubsection{$Q$ 的跨墙}
我们现在来证明一条当 $Q$ 在 $M_\R$ 中变动时的跨墙公式。为此我们需要理解以 $Q$ 为终点的断线如何随 $Q$ 变动。
为此我们引入断线的形变与退化断线的概念：

\begin{definition}
一\emph{族断线}由如下数据组成：
\begin{itemize}[leftmargin=*]
\item 连续映射 $B : (-\infty,0]\times I \to M_\R$，其中 $I\subseteq\R$ 是区间。
\item 连续函数 $t_0,\dots,t_n : I\to[-\infty,0]$，使得对 $s\in I$ 有
$-\infty = t_0(s)<t_1(s)<\dots<t_n(s) = 0$。
\item 单项式 $c_iz^{m^B_i}$，$1\le i\le n$。
\end{itemize}
这些数据满足：对一切 $s\in I$，$B_s := B|_{(-\infty,0]\times\{s\}}$ 是通常意义下的断线，其数据为
$t_0(s)<\dots<t_n(s)$ 与单项式 $c_iz^{m^B_i}$，$1\le i\le n$。

对 $s,s'\in I$，称 $B_{s'}$ 是 $B_s$ 的一个\emph{形变}。
\end{definition}

\begin{definition}
\emph{退化断线}是在 $\operatorname{Sing}(\mathfrak D(\Sigma,P_1,\dots,P_k))$ 的某点处折弯的断线之极限。更精确
地说，退化断线是连续逆紧映射 $\beta : (-\infty,0]\to M_\R$ 连同数据 $-\infty = t_0<t_1\le\dots\le t_n=0$ 与
单项式 $c_iz^{m^\beta_i}$，使得存在：
\begin{itemize}[leftmargin=*]
\item 连续映射 $B : (-\infty,0]\times[0,1]\to M_\R$；
\item 连续函数 $\bar t_0,\dots,\bar t_n : [0,1]\to[-\infty,0]$，使得
\[
-\infty = \bar t_0(s)\le\bar t_1(s)\le\dots\le\bar t_n(s) = 0
\]
% [原文 p.208]
对 $s\in[0,1]$ 成立，且当 $s<1$ 时为严格不等式。
\end{itemize}
这些数据满足：
\begin{itemize}[leftmargin=*]
\item $B_s := B|_{(-\infty,0]\times\{s\}}$（带数据 $c_iz^{m^\beta_i}$）对 $s<1$ 是通常意义下的断线；
\item $\beta = B_1$ 且 $\bar t_i(1) = t_i$ 对一切 $i$；
\item 对某个 $i$ 有 $\beta(t_i)\in\operatorname{Sing}(\mathfrak D(\Sigma,P_1,\dots,P_k))$。
\end{itemize}
\end{definition}

注意在取这样的极限时，各个 $i$ 的 $\bar t_{i-1}$ 与 $\bar t_i$ 可能重合，因此极限可能只有更少的线性段。

现在我们有了跨墙定理：

\begin{theorem}
若 $Q,Q'\in M_\R\setminus\operatorname{Supp}(\mathfrak D(\Sigma,P_1,\dots,P_k))$ 是一般点，且 $\gamma$ 是连接
$Q$ 与 $Q'$ 的、使 $\theta_{\gamma,\mathfrak D(\Sigma,P_1,\dots,P_k)}$ 有定义的路径，则
\[
W_k(Q') = \theta_{\gamma,\mathfrak D(\Sigma,P_1,\dots,P_k)}(W_k(Q)).
\]
\end{theorem}

\begin{proof}
令 $\mathfrak D = \mathfrak D(\Sigma,P_1,\dots,P_k)$。令 $\mathfrak U$ 为 $\operatorname{Supp}(\mathfrak D)$ 与所有退化断线（端点任意）的像之并的并集。显然 $\dim\mathfrak U\le1$（当然，当 $k\ge1$ 时等于 $1$）。

注意断线 $\beta$ 总可以连续形变。这可以如下做到。我们平移 $\beta$ 的初始射线
$\beta((-\infty,t_1])$。归纳地，这就形变了 $\beta$ 的其余各段。只要某一个折弯点不碰到 $\mathfrak D$ 的奇点，
每个折弯点就始终落在 $\mathfrak D$ 中完全相同的一组射线里，从而形变后的断线可以完全照 $\beta$ 的方式折弯。
因此，只有当 $\beta$ 的这个形变收敛到一条在 $\operatorname{Sing}(\mathfrak D)$ 的奇点处折弯的退化断线时，我们
才会遇到麻烦——因为那时含折弯点的射线集合可能跳跃。

由此可知，只要 $\beta$ 的终点停留于 $M_\R\setminus\mathfrak U$ 的同一个连通分支内，$\beta$ 就可以连续形变。
更精确地说，若考虑路径 $\gamma : [0,1]\to u$（$u$ 是 $M_\R\setminus\mathfrak U$ 的连通分支），且 $\beta$ 是
以 $\gamma(0)$ 为终点的断线，则存在连续形变 $B$，使得 $B_0 = \beta$ 且 $B_s(0) = \gamma(s)$，$0\le s\le1$。

由命题 5.32，对一般的 $Q$，以 $Q$ 为边界的 Maslov 指数二圆盘与以 $Q$ 为终点的断线一一对应。因此由上面的讨论，
$W_k(Q)$ 在 $M_\R\setminus\mathfrak U$ 的同一连通分支内对一般的 $Q$ 为常值。

下面我们仔细分析当断线的终点在 $M_\R\setminus\mathfrak U$ 的不同连通分支之间穿行时，断线如何变化。于是考虑
$M_\R\setminus\mathfrak U$ 的两个连通分支 $u_1$ 与 $u_2$。令 $L = \overline{u_1}\cap\overline{u_2}$，并设
$\dim L = 1$。设 $Q_1$ 与 $Q_2$ 分别是 $u_1$ 与 $u_2$ 中接近 $L$ 的一般点，位于 $L$ 的两侧。设
$\gamma : [0,1]\to M_\R$ 是连接 $Q_1$ 与 $Q_2$、恰好穿过 $L$ 一次的短的一般路径。设 $s_0$ 是唯一使
$\gamma(s_0)\in L$ 的时刻。通过把 $\gamma$ 取得足够一般，可以假定 $\gamma(s_0)$ 是一个点，在其邻域内
$\mathfrak U$ 是流形。

设 $\cB(Q_i)$ 是以 $Q_i$ 为终点的断线之集。取 $n_0\in N$ 为湮灭 $L$ 的切空间、且在 $Q_1$ 处取值小于在 $Q_2$
处的本原向量。我们可以把 $\cB(Q_i)$ 分解为三个集合 $\cB^+(Q_i)$、$\cB^-(Q_i)$ 与 $\cB^0(Q_i)$ 如下。对
$\beta\in\cB(Q_i)$，令 $m_\beta = \beta_*(-\partial/\partial t|_{t=0})$。则依
$\langle n_0,m_\beta\rangle>0$、$<0$、$=0$ 分别有 $\beta\in\cB^+(Q_i)$、$\cB^-(Q_i)$、$\cB^0(Q_i)$。这

% [原文 p.209]
给出分解
\begin{align*}
W_k(Q_1) &= W^-_k(Q_1)+W^0_k(Q_1)+W^+_k(Q_1)\\
W_k(Q_2) &= W^-_k(Q_2)+W^0_k(Q_2)+W^+_k(Q_2).
\end{align*}
我们将证明
\begin{equation}
\theta_{\gamma,\mathfrak D}\bigl(W^-_k(Q_1)\bigr) = W^-_k(Q_2), \tag{5.10}
\end{equation}
\begin{equation}
\theta^{-1}_{\gamma,\mathfrak D}\bigl(W^+_k(Q_2)\bigr) = W^+_k(Q_1), \tag{5.11}
\end{equation}
\begin{equation}
W^0_k(Q_2) = W^0_k(Q_1). \tag{5.12}
\end{equation}
% [Tutor Note: 原文 p.209，(5.10)–(5.12)。三条的分工：$\cB^-$（$\langle n_0,m_\beta\rangle<0$）整族沿 $\gamma$
%  连续搬到 $L$ 另一侧，故 $\theta_{\gamma,\mathfrak D}$ 把 $W^-_k(Q_1)$ 送到 $W^-_k(Q_2)$；$\cB^+$ 同理反向；
%  而 $\cB^0$（$\langle n_0,m_\beta\rangle=0$，断线沿 $L$ 方向出射）跨过 $L$ 时断线会跳跃，无法整族搬过去，
%  只能逐对配对，这就是 (5.12) 必须单独证明的原因（见 p.214 的分拆）。核对锚点：
%  $m_\beta=\beta_*(-\partial/\partial t|_{t=0})$ 是断线在终点处的出射方向，故 $\langle n_0,m_\beta\rangle$ 的
%  符号判断它落在 $L$ 的哪一侧。]

由此即得所要的恒等式
\[
\theta_{\gamma,\mathfrak D}\bigl(W_k(Q_1)\bigr) = W_k(Q_2),
\]
因为 $\theta_{\gamma,\mathfrak D}$ 在 $W^0_k(Q_1)$ 上必为恒等。然后归纳地使用这一点，即可看出它对端点落在
$M_\R\setminus\mathfrak U$ 中、且使 $\theta_{\gamma,\mathfrak D}$ 有定义的任何路径 $\gamma$ 都成立。

\para{式 (5.10) 与 (5.11) 的证明。}若 $\beta$ 是以 $Q_1$ 为终点的断线，则当
$\langle n_0,m_\beta\rangle\le0$ 时 $\beta([t_{n-1},0])\cap L = \emptyset$，而当
$\langle n_0,m_\beta\rangle>0$ 时 $\beta([t_{n-1},0])\cap L\ne\emptyset$。（这里我们用到了 $Q_1$ 非常接近
$L$。）另一方面，若 $\beta$ 的终点为 $Q_2$，则当 $\langle n_0,m_\beta\rangle\ge0$ 时
$\beta([t_{n-1},0])\cap L = \emptyset$，而当 $\langle n_0,m_\beta\rangle<0$ 时
$\beta([t_{n-1},0])\cap L\ne\emptyset$。

为证（比如说）(5.10)，我们如下进行。设 $\beta\in\cB^-(Q_1)$。由上一段，
$\beta([t_{n-1},0])\cap L = \emptyset$。设 $c_nz^{m^\beta_n}$ 是联系于 $\beta$ 最后一段的单项式，并把
$\theta_{\gamma,\mathfrak D}(c_nz^{m^\beta_n})$ 按定义 5.30 (4) 写成单项式之和 $\sum_{i=1}^s d_iz^{m_i}$。于是
我们可以沿 $\gamma$ 把 $\beta$ 连续形变到时刻 $s_0$。事实上，由 $\mathfrak U$ 的定义，若 $\beta$ 收敛到一条
穿过 $\operatorname{Sing}(\mathfrak D)$ 的退化断线，则该断线的像会含于 $\mathfrak U$ 中，于是已经包含 $L$
的 $\mathfrak U$ 在 $\gamma(s_0)$ 的邻域内就不是流形了。

设 $\beta'$ 是终点为 $\gamma(s_0)$ 的 $\beta$ 的形变。对 $1\le i\le s$，我们通过沿方向 $-r(m_i)$ 给 $\beta'$
添加一条短线段、并附带单项式 $d_iz^{m_i}$，得到一条断线 $\beta'_i$。这条新断线的终点在 $u_2$ 中，因此可以形变
为断线 $\beta''_i\in\cB^-(Q_2)$。我们指出：若 $d_iz^{m_i}$ 正是 $\theta_{\gamma,\mathfrak D}(c_nz^{m^\beta_n})$
中出现的那一项 $c_nz^{m^\beta_n}$，则该断线未必真的在 $L$ 处折弯。见图 13。

反之，任何断线 $\beta\in\cB^-(Q_2)$ 显然都以这种方式出现。

由此 (5.10) 便清楚了。(5.11) 完全相同。

\para{式 (5.12) 的证明。}我们将证明存在分拆
$\cB^0(Q_1) = \coprod_{i=1}^s\cB^i_1$ 与 $\cB^0(Q_2) = \coprod_{i=1}^s\cB^i_2$，使得对每个 $i$，由
$\cB^i_1$ 与 $\cB^i_2$ 对 $W_k(Q_1)$ 与 $W_k(Q_2)$ 的贡献相同。

为叙述简单起见，我们只在“任何以 $\gamma(s_0)$ 为终点的退化断线至多经过
$\operatorname{Sing}(\mathfrak D)$ 的一个点”这一情形下描述；一般情形的处理留给读者：那只是在记号上更复杂，
概念上并不更难。

设 $\beta_1\in\cB^0(Q_1)$。若 $\beta_1$ 可连续形变为 $\cB^0(Q_2)$ 中的断线 $\beta_2$，则 $\beta_1$ 与
$\beta_2$ 各自在分拆中作为单元集出现，比如说 $\beta_1\in\cB^i_1$、$\beta_2\in\cB^i_2$，而显然这两个集合对
$W_k(Q_1)$ 与 $W_k(Q_2)$ 贡献相同的项。

现在设 $\beta_1\in\cB^0(Q_1)$ 不能连续形变为 $\cB^0(Q_2)$ 中的任何 $\beta_2$。这意味着 $\beta_1$ 必在时刻
$s_0$ 形变为一条退化断线：即存在如定义 5.34 中的
$B : (-\infty,0]\times[0,s_0]\to M_\R$，使得 $B|_{(-\infty,0]\times[0,s_0)}$ 是连续形变，而 $B_{s_0}$ 是
一条退化断线
% [原文 p.210]
在 $\operatorname{Sing}(\mathfrak D)$ 的点 $P$ 处折弯。因此必有某个 $j$ 使
$B(\bar t_j(s_0),s_0) = P\in\operatorname{Sing}(\mathfrak D)$。

我们需要分析两种情形：或者 $P\in\{P_1,\dots,P_k\}$，或者 $P\notin\{P_1,\dots,P_k\}$。

\begin{figure}[htbp]
\centering
\includegraphics[width=0.86\textwidth]{f5_13.png}
\figtitle{13}{以 $Q_1$ 与 $Q_2$ 为终点的断线。}
\end{figure}

\para{情形 1。}$P = P_i$ 对某个 $i$ 成立。由于以 $P_i$ 为端点的所有射线都涉及单项式 $u_i$，一条断线至多沿
一条以 $P_i$ 为端点的射线折弯；而如上所观察，$\beta_1$ 至少沿一条这样的射线折弯。于是把 $\beta_1$ 所沿折弯的
那条以 $P_i$ 为端点的射线记作 $\mathfrak d_1$，它对应于一棵 Maslov 指数零树
$\tilde h_1 : \tilde\Gamma_1\to M_\R$。这棵树经过 $P_i$；把它在 $P_i$ 处切开并去掉映到 $P_i$ 的标记边，
就得到以 $P_i$ 为边界的 Maslov 指数二圆盘 $h_1 : \Gamma'_1\to M_\R$。设 $\bar\beta_2$ 是对应于这个 Maslov
指数二圆盘、以 $P_i$ 为终点的断线。见图 14。

接着，回忆 $B(\bar t_j(s_0),s_0) = P_i$；令 $\bar\beta'_2 : [\bar t_j(s_0),0]\to M_\R$ 是 $B$ 在
$[\bar t_j(s_0),0]\times\{s_0\}$ 上的限制：这是从 $P_i$ 出发的一段断线。于是我们可以把 $\bar\beta_2$ 与
$\bar\beta'_2$ 拼接起来（把 $\bar\beta_2$ 定义域中的 $0$ 与 $\bar\beta'_2$ 定义域中的 $\bar t_j(s_0)$ 等同），
得到我们所希望的、经过 $P_i$ 的退化断线 $\beta'_2$。

注意 $B|_{(-\infty,\bar t_j(s_0)]\times\{s_0\}}$ 是以 $P_i$ 为终点的断线，因此对应于以 $P_i$ 为边界的
Maslov 指数二圆盘 $h_2 : \Gamma'_2\to M_\R$。把 $\Gamma'_2$ 的边 $E'_{\mathrm{out}}$ 延拓为无界边，得到热带树
$\tilde h_2 : \tilde\Gamma_2\to M_\R$；在 $\tilde\Gamma_2$ 上标出映到 $P_i$ 的点后，它成为 Maslov 指数零树，
从而对应于 $\mathfrak D$ 中以 $P_i$ 为端点的一条射线 $\mathfrak d_2$。

注意联系于 $\mathfrak d_i$ 的函数是 $1+w_{\Gamma_i}(E_{\mathrm{out},i})u_i\operatorname{Mono}(h_i)$。另一方面，
联系于 $B|_{(-\infty,\bar t_j(s_0)]\times\{s_0\}}$ 最后一段的单项式（即 $c_jz^{m^{\beta_1}_j}$）是
$\operatorname{Mono}(h_2)$，而联系于 $\bar\beta_2$ 最后一段的单项式是 $\operatorname{Mono}(h_1)$。因此特别地，
单项式 $c_{j+1}z^{m^{\beta_1}_{j+1}}$ 由 $\beta_1$ 在 $\mathfrak d_1$ 处的折弯得到，于是等于
\begin{equation}
w_{\Gamma_1}(E_{\mathrm{out},1})\langle n_1,r(\Delta(h_2))\rangle u_i\operatorname{Mono}(h_1)\operatorname{Mono}(h_2). \tag{5.13}
\end{equation}

% [原文 p.211]
\begin{figure}[htbp]
\centering
\includegraphics[width=0.80\textwidth]{f5_14.png}
\figtitle{14}{断线经过奇点 $P\in\{P_1,\dots,P_k\}$ 时的变化。}
\end{figure}

这里 $n_1\in N$ 本原、与 $\mathfrak d_1$ 正交，且在 $r(\Delta(h_2))$ 上为正。

现在我们可以如下形变 $\beta'_2$：把 $\bar\beta_2$ 的端点沿 $\mathfrak d_2$ 移离 $P_i$，并带动 $\bar\beta'_2$
一起移动。不过我们还需要追踪单项式：必须确保 $\bar\beta'_2$ 第一段上的单项式正是 $\beta'_2$ 沿 $\mathfrak d_2$
折弯时会出现的那个。然而后一个单项式是把联系于穿过 $\mathfrak d_2$ 的自同构作用于 $\operatorname{Mono}(h_1)$
所得的一项，故正好是
\begin{equation}
w_{\Gamma_2}(E_{\mathrm{out},2})\langle n_2,r(\Delta(h_1))\rangle u_i\operatorname{Mono}(h_1)\operatorname{Mono}(h_2). \tag{5.14}
\end{equation}
同样，$n_2\in N$ 本原、与 $\mathfrak d_2$ 正交，且在 $r(\Delta(h_1))$ 上为正。事实上，对 $i=1,2$，
\[
w_{\Gamma_i}(E_{\mathrm{out},i})n_i = \pm X_{r(\Delta(h_i))},
\]
符号的选取使 (5.13) 与 (5.14) 都等于
\[
|r(\Delta(h_1))\wedge r(\Delta(h_2))|u_i\operatorname{Mono}(h_1)\operatorname{Mono}(h_2).
\]
于是 $\beta'_2$ 现在可以形变离开奇点 $P_i$，得到以 $Q_2$ 为终点的断线 $\beta_2$。注意这绝不表示连续形变：
断线在经过 $P_i$ 时确实发生了跳跃。

注意这一过程是可逆的。若从 $\beta_2$ 出发、试图如上把它形变穿过 $P_i$，就得到 $\beta_1$。

总之，在这一情形我们可以取形如 $\beta_1\in\cB^i_1$ 与 $\beta_2\in\cB^i_2$（对某个 $i$）的单元集作为分拆。
它们对 $W_k(Q)$ 给出相同的贡献。

\para{情形 2。}$P\notin\{P_1,\dots,P_k\}$。设 $\mathfrak D_P$ 是命题 5.28 的证明中所构造的散射图，它由以下
元素组成：对 $\mathfrak D(\Sigma,P_1,\dots,P_k)$ 中每条以 $P$ 为初始点的射线，取一条以原点为初始点的射线；
对 $\mathfrak D(\Sigma,P_1,\dots,P_k)$ 中每条经过 $P$、且初始点不等于 $P$ 的射线，取一条过原点的直线。

% [原文 p.212]
由命题 5.28，
\begin{equation}
\theta_{\gamma,\mathfrak D_P} = \mathrm{Id} \tag{5.15}
\end{equation}
其中 $\gamma$ 是绕 $P$ 的环路。

由于 $\mathfrak D_P$ 的行为比较简单，可以通过仔细的逐例分析来分析这一情形：这已在 [42] 中完成。这里我们给出
一个更简单、更一般的途径，它来自 Carl、Pumperla 与 Siebert 正在进行的工作 [13]。

把 $M_\R$ 中的\emph{轨道}定义为一对 $(\mathfrak t, cz^m)$，其中 $cz^m\in\C[T_\Sigma]\otimes R_k$ 是单项式，
且或者
\begin{enumerate}[label=(\arabic*),leftmargin=*]
\item $\mathfrak t = \R_{\ge0}r(m)$，此时称 $(\mathfrak t,cz^m)$ 为\emph{入射轨道}，或者
\item $\mathfrak t = -\R_{\ge0}r(m)$，此时称 $(\mathfrak t,cz^m)$ 为\emph{出射轨道}。
\end{enumerate}
照例，$\mathfrak t$ 称为 $(\mathfrak t,cz^m)$ 的\emph{支撑}。

\begin{claim}
给定一条入射轨道 $\mathfrak t_0 = (\R_{\ge0}r(m_0),c_0z^{m_0})$，存在一族出射轨道
$\{\mathfrak t_i\}$，其中 $\mathfrak t_i = (-\R_{\ge0}r(m_i),c_iz^{m_i})$，具有如下性质。对每条 $\mathfrak t_i$
（包括 $i=0$），设 $x_i$ 是 $M_\R\setminus\operatorname{Supp}(\mathfrak D_P)$ 的某个连通分支中的点，该分支的闭包
含有 $\mathfrak t_i$ 的支撑。设 $\gamma_i$ 是从 $x_i$ 到 $x_0$、不经过 $P$ 的路径；由 (5.15)，
$\theta_i := \theta_{\gamma_i,\mathfrak D_P}$ 与此路径的选取无关。则
\begin{equation}
c_0z^{m_0} = \sum_i\theta_i(c_iz^{m_i}). \tag{5.16}
\end{equation}
此外，这族轨道是唯一的：给定 $M_\R$ 中的射线 $\R_{\ge0}m$，在对一切 $i\ne0$、$\mathfrak t_i$ 的支撑为
$\R_{\ge0}m$ 的下标求和所得的和 $\sum_i c_iz^{m_i}$ 由 $\mathfrak t_0$ 唯一确定。
\end{claim}

\begin{proof}
首先注意 (5.16) 不依赖于 $x_i$ 或 $x_0$ 的精确选取：若 $\mathfrak t_i$（包括 $i=0$ 的情形）含于 $\mathfrak D_P$
的一条射线或直线中，则 $c_iz^{m_i}$ 在相应的自同构下不变。

我们现归纳地构造集合 $\{\mathfrak t_i\}$。第 $p$ 步我们将找到一族出射轨道 $\fT_p$，使得 (5.16) 在模理想
$\mathfrak m^p$ 的意义下成立，其中 $\mathfrak m$ 是由 $u_1,\dots,u_k$ 生成的理想。对 $p=1$，取
\[
\fT_1 = \{(-\R_{\ge0}r(m_0),c_0z^{m_0})\},
\]
这是可行的，因为 $\theta_i\equiv\mathrm{Id}\mod\mathfrak m$。

今设已构造出 $\fT_p = \{\mathfrak t_i \mid i\in I\}$。则由归纳假设，
\[
c_0z^{m_0}-\sum_i\theta_i(c_iz^{m_i}) = \sum c_jz^{m_j} \mod \mathfrak m^{p+1},
\]
其中 $c_j\in\mathfrak m^p$。于是取
\[
\fT_{p+1} = \fT_p\cup\{(-\R_{\ge0}r(m_j),c_jz^{m_j})\};
\]
此时 (5.16) 模 $\mathfrak m^{p+1}$ 成立。

% [Tutor Note: 原文 p.212，断言 5.36 的证明。这里用到 $\mathfrak m=(u_1,\dots,u_k)$ 的幂零性：由 $u_i$ 平方为零，
%  任意 $k+1$ 个生成元的乘积为零，故 $\mathfrak m^{k+1}=0$，$\mathfrak m$-进归纳在 $k+1$ 步后必然终止；
%  “每一步都没有可作的选择”给出唯一性。核对锚点：第 $p$ 步是把余项
%  $c_0z^{m_0}-\sum_i\theta_i(c_iz^{m_i})$ 中 $\mathfrak m^p$ 阶的项逐项吸收为新轨道，步数受 $\mathfrak m^{k+1}=0$ 限制。]

由于 $\mathfrak m^{k+1}=0$，这个过程会终止。同样清楚的是，每一步都没有可作的选择，因此唯一性成立。
\end{proof}

这族出射轨道的意义在于，它精确地告诉我们：靠近 $P$ 的断线在 $P$ 附近折弯时能产出哪些项。更精确地说，

% [原文 p.213]
\begin{figure}[htbp]
\centering
\includegraphics[width=0.62\textwidth]{f5_15.png}
\figtitle{15}{给定初始单项式的局部断线之集 $\cB$。}
\end{figure}

我们考虑一种局部形式的断线：它们是像与 $\operatorname{Sing}(\mathfrak D_P)$ 不相交的连续映射
$\beta : \R\to M_\R$，连同数据
\[
-\infty = t_0<\dots<t_n=\infty
\]
以及与定义 5.30 中相同的单项式 $c_iz^{m^\beta_i}$，并满足该定义的条件 (1)、(3) 与 (4)。换言之，我们不限制
初始单项式，也没有终点。除此之外，$\beta$ 就是一条断线。我们称之为\emph{局部断线}。

固定单项式 $c_0z^{m_0}$ 与一般点 $x\in M_\R$，其中 $x$ 含于 $M_\R\setminus\operatorname{Sing}(\mathfrak D_P)$ 的某个连通分支中，且该分支的闭包含有 $\R_{\ge0}r(m_0)$。考虑（必然有限的）集合 $\cB$，它由这样的局部断线 $\beta$
组成：联系于第一段 $(t_0,t_1)$ 的单项式为 $c_0z^{m_0}$，且存在 $t\in(t_0,t_1)$ 使 $\beta(t) = x$。
（若两条局部断线只差一个时间平移，我们就视它们为同一条。）见图 15。对 $\beta_i\in\cB$，设 $c_iz^{m_i}$ 是
联系于 $\beta_i$ 最后一段线段的单项式，并设 $\gamma_i$ 是从点 $x_i = \beta_i(t)$（$t\gg0$）到 $x$ 的路径。令
$\theta_i = \theta_{\gamma_i,\mathfrak D_P}$。

\begin{claim}
$\sum_i\theta_i(c_iz^{m_i}) = c_0z^{m_0}$。
\end{claim}

\begin{proof}
注意若 $\mathfrak D_P$ 只由一条直线组成，则 $\cB$ 中的断线至多在一个点处折弯，即 $x-\R_{\ge0}r(m_0)$ 与
$\mathfrak D_P$ 唯一元素的交点。此时由定义 5.30 的条件 (4) 立刻可知断言成立。

利用这一点，一般情形就容易得到：只需对路径 $\gamma_i$ 取 $\beta_i$ 从 $x_i$ 倒退到 $x$ 的那一段。
\end{proof}

固定一个出射方向 $-\R_{\ge0}m$，并考虑和 $\sum_i c_iz^{m_i}$，其中 $i$ 取遍所有满足
$-r(m_i)\in-\R_{\ge0}m$ 的下标。作为断言 5.36 唯一性陈述的推论，这个和是良定的，与点 $x$ 的选取无关。

现在我们可以完成这一情形的证明。回到 (5.12) 证明开头的局面：我们有 $\beta_1\in\cB^0(Q_1)$，它形变为一条
经过 $P$ 的退化断线 $B_{s_0}$。

% [原文 p.214]
我们取 $\cB^i_1\subseteq\cB^0(Q_1)$ 为 $\cB^0(Q_1)$ 中所有形变到与 $B_{s_0}$ 相同的极限映射的断线之集（但可以
涉及不同的单项式）。类似地定义 $\cB^i_2$ 为 $\cB^0(Q_2)$ 中所有形变到同一极限映射的断线之集。这两个集合分别对
$W_k(Q_1)$ 与 $W_k(Q_2)$ 给出贡献；一旦证明这两个贡献相同，就完成了。

回忆我们沿路径 $\gamma$ 形变 $\cB^0(Q_1)$ 与 $\cB^0(Q_2)$ 中断线的端点；于是对 $s\ne s_0$，定义 $\cB_s$ 为
终点为 $\gamma(s)$、且同样形变到极限映射 $B_{s_0}$ 的断线之集。这样，$s<s_0$ 时 $\cB_s$ 中每条断线都是
$\cB^i_1$ 中唯一一条断线的形变，而 $s>s_0$ 时 $\cB_s$ 中每条断线都是 $\cB^i_2$ 中唯一一条断线的形变。

令 $x_0 := B_{s_0}(t)$，其中 $t$ 比使 $B_{s_0}(\bar t) = P$ 的那个 $\bar t$ 略小。在 $x_0$ 附近局部地，
$B_{s_0}$ 的像把 $M_\R$ 分成两个连通分支；$s<s_0$ 时 $\cB_s$ 中的断线局部含于其中一个分支，$s>s_0$ 时
$\cB_s$ 中的断线局部含于另一个分支。取 $x_0$ 附近 $B_{s_0}$ 的像两侧的点 $x$ 与 $x'$，使得对某个 $s<s_0$，
$\cB_s$ 的某个元素经过（比如说）$x$。对每条断线 $\beta\in\cB^i_1$，可以找到至少一个 $s<s_0$，使得 $\beta$
能形变为经过 $x$ 的某条 $\beta_s\in\cB_s$。类似地，对每条断线 $\beta\in\cB^i_2$，可以找到至少一个 $s>s_0$，
使得 $\beta$ 能形变为经过 $x'$ 的某条 $\beta_s\in\cB_s$。对每条 $\beta\in\cB^i_1$，我们在这些经过 $x$ 的
$\beta_s$ 中作一个选取，并令 $\cB_x$ 为这些选取之集。类似地定义 $\cB_{x'}$，使得 $\cB_x$ 与 $\cB^i_1$ 一一对应、
$\cB_{x'}$ 与 $\cB^i_2$ 一一对应。特别地，若 $c_\beta z^{m_\beta}$ 表示联系于断线 $\beta$ 最后一段的单项式，
我们只需证明
\[
\sum_{\beta\in\cB_x}c_\beta z^{m_\beta} = \sum_{\beta\in\cB_{x'}}c_\beta z^{m_\beta}.
\]
为证这一点，只需证明把 $c_\beta z^{m_\beta}$ 换成联系于 $\beta$ 在 $P$ 附近折弯完毕之后那一刻的单项式后，同样
的陈述成立。结果现在由断言 5.37 与断言 5.36 的唯一性陈述推出。
\end{proof}

\subsubsection{变动 $P_1,\ldots,P_k$}
现在我们转而研究 $W_k(Q)$ 对点 $P_1,\dots,P_k$ 的依赖关系。为此我们需要散射图的三维版本。

\begin{definition}
设 $L\subseteq\R$ 是闭区间。设 $\pi_1$ 与 $\pi_2$ 是 $M_\R\times L$ 到 $M_\R$ 与 $L$ 的投影。
$M_\R\times L$ 中的\emph{散射图}是一个有限集 $\mathfrak D$，由对 $(\mathfrak d,f_{\mathfrak d})$ 组成，使得
\begin{itemize}[leftmargin=*]
\item $\mathfrak d\subseteq M_\R\times L$ 是二维多面体，且 $\pi_2(\mathfrak d)$ 是一维的。此外存在一维子集
$\mathfrak b\subseteq M_\R\times L$ 与元素 $m_0\in T_\Sigma$（$r(m_0)\ne0$），使得
\[
\mathfrak d = \mathfrak b-\R_{\ge0}(r(m_0),0).
\]
\item $f_{\mathfrak d}\in\C[z^{m_0}]\otimes_\C R_k\subseteq\C[T_\Sigma]\otimes_\C R_k[\![y_0]\!]$。
\item $f_{\mathfrak d}\equiv1\mod(u_1,\dots,u_k)z^{m_0}$。
\end{itemize}
我们定义
\[
\operatorname{Sing}(\mathfrak D) = \bigcup_{\mathfrak d\in\mathfrak D}\partial\mathfrak d
\ \cup\bigcup_{\substack{\mathfrak d_1,\mathfrak d_2\in\mathfrak D\\ \dim\mathfrak d_1\cap\mathfrak d_2=1}}
\mathfrak d_1\cap\mathfrak d_2.
\]
这是 $M_\R\times L$ 的一维子集。设 $\operatorname{Interstices}(\mathfrak D)$ 为使
$\operatorname{Sing}(\mathfrak D)$ 不是流形的有限个点构成之集。我们用
$\operatorname{Joints}(\mathfrak D)$ 记

% [原文 p.215]
$\operatorname{Sing}(\mathfrak D)\setminus\operatorname{Interstices}(\mathfrak D)$ 的连通分支的闭包之集，
并把 $\operatorname{Joints}(\mathfrak D)$ 与 $\operatorname{Interstices}(\mathfrak D)$ 的元素分别称为
\emph{接缝}与\emph{间隙}。\footnote{这一术语取自 [49]。}称一条接缝是\emph{水平的}，如果它在 $\pi_2$ 下的像是
一个点；否则称该接缝是\emph{竖直的}。
\end{definition}

% [Tutor Note: 原文 p.214–215，定义 5.38。三维散射图与二维的唯一差别是：$\mathfrak d$ 换成二维多面体、
%  $\pi_2(\mathfrak d)$ 一维，且 $\mathfrak d=\mathfrak b-\R_{\ge0}(r(m_0),0)$——所以每面墙都沿
%  $(\R_{\ge0}(r(m_0),0))$ 方向延展，即沿 $L$ 方向的“竖直墙”不会出现。这正是定理 5.39 证明中“$\widetilde{\mathfrak D}$
%  中没有投影到 $L$ 中点的墙”的来源。核对锚点：$-\R_{\ge0}$ 而非 $+\R_{\ge0}$ 使 $\mathfrak d$ 成为“半平面带”，
%  从而 $\pi_2(\mathfrak d)$ 是一维而非二维。]

对 $(M_\R\times L)\setminus\operatorname{Sing}(\mathfrak D)$ 中的路径 $\gamma$，可以完全像 $M_\R$ 中散射图
的情形那样定义一个元素 $\theta_{\gamma,\mathfrak D}\in \mathbb V_{\Sigma,k}$。事实上，我们只需定义 $\gamma$ 在时刻
$t_i$ 穿过 $(\mathfrak d,f_{\mathfrak d})$ 时的自同构 $\theta_{\gamma,\mathfrak d}$。设 $\gamma$ 从 $\mathfrak d$
的一侧穿到另一侧，照例取本原的 $n_0\in N$，使 $\langle n_0,r(m_0)\rangle=0$，且 $n_0$ 在 $\gamma$ 穿入的那
一侧 $\mathfrak d$ 上取值较小。于是照例还可以定义
\[
\theta_{\gamma,\mathfrak d}(z^m) = z^mf_{\mathfrak d}^{\langle n_0,r(m)\rangle}.
\]
同样容易验证，$\theta_{\gamma,\mathfrak D}$ 只依赖于路径 $\gamma$ 在
$(M_\R\times L)\setminus\operatorname{Sing}(\mathfrak D)$ 中的同伦型。

$M_\R\times L$ 中的\emph{断线}是映射 $\beta : (-\infty,0]\to M_\R\times L$ 连同数据 $t_0<\dots<t_n$ 与
单项式 $c_iz^{m^\beta_i}$，使得
\begin{enumerate}[label=(\arabic*),leftmargin=*]
\item $\pi_2\circ\beta$ 为常值，比如说像为 $P\in L$。
\item 在把 $M_\R\times\{P\}$ 与 $M_\R$ 等同之后，$\pi_1\circ\beta$ 是关于 $M_\R$ 中的散射图 $\mathfrak D_P$（由
\[
\mathfrak D_P := \bigl\{(\mathfrak d\cap(M_\R\times\{P\}),f_{\mathfrak d}) \bigm| (\mathfrak d,f_{\mathfrak d})\in\mathfrak D
 \text{ 且 } \mathfrak d\cap(M_\R\times\{P\})\ne\emptyset\bigr\}
\]
给出）按定义 5.30 的意义的断线。
\end{enumerate}

现在我们可以描述 $W_k(Q)$ 随点 $P_1,\dots,P_k$ 变动的方式。

\begin{theorem}
设 $W$ 与 $W'$ 分别是两组不同的一般点 $P_1,\dots,P_k$ 与 $P'_1,\dots,P'_k$ 所对应的 $W_k(Q)$。则
\[
W' = \theta(W)
\]
对某个 $\theta\in \mathbb V_{\Sigma,k}$ 成立。
\end{theorem}

\begin{proof}
我们对 $k$ 作归纳来证明这一结果。注意 $k=1$ 的初始情形是显然的，因为移动 $P_1$ 而保持 $Q$ 不动，与移动 $Q$
而保持 $P_1$ 不动是同一件事。

显然只需在只有 $P_1$ 变动的情形证明这一结果，因为然后我们可以对 $P_2,\dots,P_k$ 依次使用同一结果。于是考虑
一组一般点 $P_1,\dots,P_k$ 与 $P'_1$。考虑连接 $P_1$ 与 $P'_1$ 的线段 $L$。对 $L$ 中除有限多个点之外的所有
点 $P$，我们可以假定 $P,P_2,\dots,P_k$ 足够一般，使得
$\operatorname{Trees}(\Sigma,P,P_2,\dots,P_k)$ 有限，且该集合的所有元素都是三价的。这就给出对 $L$ 中一般的
$P$ 的一族散射图 $\mathfrak D(\Sigma,P,P_2,\dots,P_k)$。我们可以把这些散射图汇总为 $M_\R\times L$ 中的一个散射图
$\widetilde{\mathfrak D} = \widetilde{\mathfrak D}(\Sigma,L,P_2,\dots,P_k)$。$\widetilde{\mathfrak D}$ 由如下
要求确定：对 $L$ 中一般的 $P$，
\[
D(\Sigma,P,P_2,\dots,P_k)
= \bigl\{(\tilde{\mathfrak d}\cap(M_\R\times\{P\}),f_{\tilde{\mathfrak d}}) \bigm|
(\tilde{\mathfrak d},f_{\tilde{\mathfrak d}})\in\widetilde{\mathfrak D}\ \text{且}\
\tilde{\mathfrak d}\cap(M_\R\times\{P\})\ne\emptyset\bigr\}.
\]

% [原文 p.216]
为了追踪 $W_k(Q)$ 对点 $P\in L$ 的依赖，我们记 $W_k(Q;P)$。我们想证明：若 $\gamma$ 是 $M_\R\times L$ 中连接
$(Q,P_1)$ 与 $(Q,P'_1)$ 的一般路径，则
\begin{equation}
W_k(Q;P'_1) = \theta_{\gamma,\widetilde{\mathfrak D}}\bigl(W_k(Q;P_1)\bigr). \tag{5.17}
\end{equation}
% [Tutor Note: 原文 p.216，(5.17)。归纳基础 $k=1$ 之所以“显然”，是因为“移动 $P_1$ 而固定 $Q$”与“移动 $Q$ 而
%  固定 $P_1$”是同一族（都在 $M_\R\times L$ 里换个方向走），故定理 5.35 直接给出。证明 (2) 时特意要求 $P=P'$，
%  于是墙是竖直的、$\pi_2$ 在墙上是常值，正好退回定理 5.35 的二维情形。]

只需证明
\begin{enumerate}[label=(\arabic*),leftmargin=*]
\item $W_k(Q;P)$ 当 $(Q,P)$ 在 $(M_\R\times L)\setminus\operatorname{Supp}(\widetilde{\mathfrak D})$ 的同一连通
分支内变动时为常值。
\item 对由一道墙 $(\tilde{\mathfrak d},f_{\tilde{\mathfrak d}})$ 分隔的两个这样的连通分支，以及墙两侧的点
$(Q,P)$、$(Q',P')$，有
\[
W_k(Q';P') = \theta_{\gamma,\mathfrak D}\bigl(W_k(Q;P)\bigr),
\]
其中 $\gamma$ 是连接 $(Q,P)$ 与 $(Q',P')$ 的短路径。
\end{enumerate}
一旦证明了 (1)，定理 5.35 就已给出 (2)：由于 $\widetilde{\mathfrak D}$ 中没有投影到 $L$ 中点的墙，我们总可以
选取位于某道墙两侧、且 $P = P'$ 的点 $(Q,P)$、$(Q',P')$，于是归结为定理 5.35 已证的情形。所以只需证明 (1)。

为证 (1)，我们用研究 $Q$ 变动时的同一技巧，即形变断线。在
$(M_\R\times L)\setminus\operatorname{Supp}(\widetilde{\mathfrak D})$ 的同一连通分支内取一般的 $(Q,P)$ 与
$(Q',P')$，并沿一般路径 $\gamma$ 从 $(Q,P)$ 走到 $(Q',P')$。考虑 $M_\R\times L$ 中以 $\gamma(t)$ 为终点的
断线。随着 $t$ 变化，我们可以连续形变以 $\gamma(t)$ 为终点的断线，除非该断线收敛到一条经过
$\widetilde{\mathfrak D}$ 的奇点的断线。然而由于这样一族断线在 $M_\R\times L$ 中描出一个二维子集，只要把
$\gamma$ 取得足够一般，就能保证这些断线中没有一个收敛到经过 $\widetilde{\mathfrak D}$ 的间隙的断线，因为间隙
的余维数是三。不过它们可以经过接缝，这需要一些小心处理。

第一个观察是：断线经过竖直接缝时会发生什么，我们已在定理 5.35 的证明中分析过。事实上，我们可以不妨假定
$\gamma$ 的选取使得当某条断线在时刻 $t_0$ 经过一条竖直接缝时，$\pi_2(\gamma(t))$ 在 $t_0$ 的邻域内保持
不变。这时我们就正好处于定理 5.35 所分析的情形。

所以我们只需考察断线经过水平接缝时会发生什么。注意水平接缝出现在散射图中两条或多条平行射线随点 $P$ 变动而
汇合之时；这通常会导致某些 $P$ 值处出现 Maslov 指数零圆盘族，或出现 Maslov 指数 $-2$ 圆盘：见图 16 与图 17。

事实上，只需证明：若 $\mathrm j$ 是一条水平接缝，而 $\gamma_{\mathrm j}$ 是 $M_\R\times L$ 中绕该接缝的小环路，则
$\theta_{\gamma_{\mathrm j},\widetilde{\mathfrak D}} = \mathrm{Id}$。的确，若 $\mathrm j$ 投影到 $P\in L$，则 $\mathrm j$ 含于某些
多边形 $\tilde{\mathfrak d}_1,\dots,\tilde{\mathfrak d}_n\in\widetilde{\mathfrak D}$ 中，且对 $P$ 附近的
$P'\in L$，$\tilde{\mathfrak d}_i\cap(M_\R\times\{P'\})$ 或者是与 $\mathrm j$ 平行的射线，或者是空集。于是当
$P'\in L$ 从 $P$ 的一侧移到另一侧时，$\mathfrak D(\Sigma,P',P_2,\dots,P_k)$ 中的某些平行射线
$\mathfrak d_1,\dots,\mathfrak d_p$ 汇合而产生该接缝，然后在 $P$ 的另一侧变成平行射线
$\mathfrak d'_1,\dots,\mathfrak d'_{p'}$。

设 $\mathfrak D_1,\mathfrak D_2$ 是 $M_\R$ 中由 $\mathfrak D(\Sigma,P',P_2,\dots,P_k)$ 给出的散射图，其中 $P'$ 非常接近 $P$，但位于
$P$ 的两侧。设 $\gamma$

% [原文 p.217]
\begin{figure}[htbp]
\centering
\includegraphics[width=0.62\textwidth]{f5_16.png}
\figtitle{16}{当 $P_1$ 向下移动时，第二幅图中突然出现一棵 Maslov 指数 $-2$ 树。此外还出现一族 Maslov 指数零树，如图 17 所示。}
\end{figure}

\begin{figure}[htbp]
\centering
\includegraphics[width=0.55\textwidth]{f5_17.png}
\figtitle{17}{}
\end{figure}

% [Tutor Note: 原文 p.217，图 16 与图 17。图 16 的四幅小图展示 $P_1$ 向下移动时散射图的变化：第二幅中两族平行
%  射线汇合，出现 Maslov 指数 $-2$ 树（即跨墙公式失效的瞬间）；第三、四幅中一族 Maslov 指数零树随之出现，
%  图 17 画出其中一条。这正是“水平接缝”处的退化：两条平行射线汇合后再分开，端点位置发生跳跃——定理 5.39 的
%  证明必须证明绕该接缝一圈回到原处时 $\theta_{\gamma_{\mathrm j},\widetilde{\mathfrak D}}=\mathrm{Id}$，否则跨墙公式不成立。]

是穿过 $\mathrm j$ 的一条短线段路径，于是可以写成
\begin{align*}
  \theta_{\gamma,\mathfrak D_1} &= \theta_{\gamma,\mathfrak d_1}\circ\dots\circ\theta_{\gamma,\mathfrak d_p},\\
  \theta_{\gamma,\mathfrak D_2} &= \theta_{\gamma,\mathfrak d'_1}\circ\dots\circ\theta_{\gamma,\mathfrak d'_{p'}}.
\end{align*}
（注意次序无关紧要，因为所有这些自同构都交换。）但
$\theta_{\gamma_{\mathrm j},\widetilde{\mathfrak D}} = \theta^{-1}_{\gamma,\mathfrak D_1}\circ\theta_{\gamma,\mathfrak D_2}$，所以若
$\theta_{\gamma_{\mathrm j},\widetilde{\mathfrak D}} = \mathrm{Id}$，就有 $\theta_{\gamma,\mathfrak D_1} = \theta_{\gamma,\mathfrak D_2}$。
这意味着，按定义 5.30 (4)，断线在 $P$ 附近、接缝 $\mathrm j$ 两侧的行为方式相同。注意两侧实际的断线集合可能不同，
因为我们并不是断言集合 $\{\theta_{\gamma,\mathfrak d_1},\dots,\theta_{\gamma,\mathfrak d_p}\}$ 与
$\{\theta_{\gamma,\mathfrak d'_1},\dots,\theta_{\gamma,\mathfrak d'_{p'}}\}$ 相同，而是断言沿这两组断线折弯的总贡献保持不变。

% 边界说明：本块止于印刷 p.217 末（跨页句按 §10 已补完："……而是断言沿这两组断线折弯的总贡献保持不变。"）。
% **重要**：定理 5.39 的证明环境在本块末尾仍未闭合（proof 已开、证明延续到 p.218 之后）。
% c05d 请从"为证明对每条水平接缝 $\theta_{\gamma_{\mathrm j},\widetilde{\mathfrak D}}=\mathrm{Id}$，我们如下进行……"
% （原文 "To show that θ_{γj,D̃} = Id for each horizontal joint, we proceed as follows."）起直接续写，
% **不要再开新的 proof 环境**，并在该定理证明结束处补上 proof 环境的结束；也不要重复本块已完成的句子。
